% GENERATED FILE. Edit canonical lessons, then run npm run generate:books.
% canonical-source-hash: fnv1a64:0548b1b08802a9b8
% canonical-lessons: GU-C426-lakh, GU-C426-karod, GU-C426-badbaki, GU-C426-gunvu, GU-C426-vadharvu

\chapter{Lakhs and Crores}
\label{ch:gu-lakhs-and-crores}
\hlchaptermodality{426}

\begin{chapteropening}
\textbf{By the end of this chapter:} \emph{I can talk about large sums and simple sums.}

Complete the last lesson of chapter 426: I can talk about large sums and simple sums.
\end{chapteropening}

\section[\texorpdfstring{\gu{લાખ} (lākh)}{લાખ (lākh)}]{\gu{લાખ} (lākh) — a lakh (one hundred thousand)}

\label{lesson:GU-C426-lakh}



\begin{quote}
Before the new one: say the Gujarati for a pledge, then the Gujarati for paid in full.
\end{quote}

\subsection*{You'll want to know: \gu{લાખ}}
\textbf{\gu{લાખ}} — \emph{lākh} — \textquotedblleft{}a lakh (one hundred thousand)\textquotedblright{}.

\begin{quote}
\gu{આ} \gu{ગાડીની} \gu{કિંમત} \gu{પાંચ} \gu{લાખ} \gu{રૂપિયા} \gu{છે}.
\end{quote}

\begin{quote}
\emph{ā gāḍīnī kimmat pā̃c lākh rūpiyā che.}
\end{quote}

\begin{quote}
This car costs five lakh rupees.
\end{quote}

\begin{cousinweb}
\textbf{\gu{લાખ}} (\emph{lākh}) is a tadbhava of Sanskrit \emph{lakṣa}, \textquotedblleft{}a mark; a hundred thousand\textquotedblright{}; CDIAL lists the Gujarati word (CDIAL 10881). Indian numbers count in lakhs and crores.
\end{cousinweb}

\begin{grammarlens}[title={What you've built}]
One more word for getting around.
\end{grammarlens}

\subsection*{Your turn}
\begin{itemize}
\raggedright
  \item \emph{Say it:} \emph{lākh}
  \item \emph{Say it:} \emph{lākh}, once more
  \item \emph{Say it:} \emph{ā gāḍīnī kimmat pā̃c lākh rūpiyā che.} — This car costs five lakh rupees.
  \item \emph{Recall:} say the Gujarati for a pledge, then the Gujarati for paid in full, then say \emph{lākh} again
\end{itemize}

\begin{tcolorbox}[breakable,colback=teal!4,colframe=teal!35!black,title={Before you move on}]
What does \gu{લાખ} mean? (\textquotedblleft{}A lakh (one hundred thousand)\textquotedblright{}.) Say it once more.
\end{tcolorbox}
\section[\texorpdfstring{\gu{કરોડ} (karoḍ)}{કરોડ (karoḍ)}]{\gu{કરોડ} (karoḍ) — a crore (ten million)}

\label{lesson:GU-C426-karod}



\begin{quote}
Before the new one: say the Gujarati for paid in full, then the Gujarati for a lakh.
\end{quote}

\subsection*{You'll want to know: \gu{કરોડ}}
\textbf{\gu{કરોડ}} — \emph{karoḍ} — \textquotedblleft{}a crore (ten million)\textquotedblright{}.

\begin{quote}
\gu{આ} \gu{શહેરમાં} \gu{એક} \gu{કરોડ} \gu{માણસો} \gu{રહે} \gu{છે}.
\end{quote}

\begin{quote}
\emph{ā śahermā̃ ek karoḍ māṇso rahe che.}
\end{quote}

\begin{quote}
Ten million people live in this city.
\end{quote}

\begin{cousinweb}
\textbf{\gu{કરોડ}} (\emph{karoḍ}) is a tadbhava of Sanskrit \emph{kōṭi}, \textquotedblleft{}ten million\textquotedblright{}; CDIAL lists the Gujarati word (CDIAL 3498). A crore is a hundred lakhs.
\end{cousinweb}

\begin{grammarlens}[title={What you've built}]
One more word for getting around.
\end{grammarlens}

\subsection*{Your turn}
\begin{itemize}
\raggedright
  \item \emph{Say it:} \emph{karoḍ}
  \item \emph{Say it:} \emph{karoḍ}, once more
  \item \emph{Say it:} \emph{ā śahermā̃ ek karoḍ māṇso rahe che.} — Ten million people live in this city.
  \item \emph{Recall:} say the Gujarati for paid in full, then the Gujarati for a lakh, then say \emph{karoḍ} again
\end{itemize}

\begin{tcolorbox}[breakable,colback=teal!4,colframe=teal!35!black,title={Before you move on}]
What does \gu{કરોડ} mean? (\textquotedblleft{}A crore (ten million)\textquotedblright{}.) Say it once more.
\end{tcolorbox}
\section[\texorpdfstring{\gu{બાદબાકી} (bādbākī)}{બાદબાકી (bādbākī)}]{\gu{બાદબાકી} (bādbākī) — subtraction}

\label{lesson:GU-C426-badbaki}



\begin{quote}
Before the new one: say the Gujarati for a lakh, then the Gujarati for a crore.
\end{quote}

\subsection*{You'll want to know: \gu{બાદબાકી}}
\textbf{\gu{બાદબાકી}} — \emph{bādbākī} — \textquotedblleft{}subtraction\textquotedblright{}.

\begin{quote}
\gu{બાળકો} \gu{શાળામાં} \gu{બાદબાકી} \gu{શીખે} \gu{છે}.
\end{quote}

\begin{quote}
\emph{bāḷko śāḷāmā̃ bādbākī śīkhe che.}
\end{quote}

\begin{quote}
Children learn subtraction at school.
\end{quote}

A feminine noun.

\begin{cousinweb}
\textbf{\gu{બાદબાકી}} (\emph{bādbākī}) is built from \gu{બાદ}, \textquotedblleft{}after, taken away\textquotedblright{}, usually traced to Arabic \emph{baʿd}, \textquotedblleft{}after\textquotedblright{}, and \gu{બાકી}, \textquotedblleft{}what remains\textquotedblright{}, usually traced to Arabic \emph{bāqī}: you take away and look at what is left.
\end{cousinweb}

\begin{grammarlens}[title={What you've built}]
One more word for getting around.
\end{grammarlens}

\subsection*{Your turn}
\begin{itemize}
\raggedright
  \item \emph{Say it:} \emph{bādbākī}
  \item \emph{Say it:} \emph{bādbākī}, once more
  \item \emph{Say it:} \emph{bāḷko śāḷāmā̃ bādbākī śīkhe che.} — Children learn subtraction at school.
  \item \emph{Recall:} say the Gujarati for a lakh, then the Gujarati for a crore, then say \emph{bādbākī} again
\end{itemize}

\begin{tcolorbox}[breakable,colback=teal!4,colframe=teal!35!black,title={Before you move on}]
What does \gu{બાદબાકી} mean? (\textquotedblleft{}Subtraction\textquotedblright{}.) Say it once more.
\end{tcolorbox}
\section[\texorpdfstring{\gu{ગુણવું} (guṇvũ)}{ગુણવું (guṇvũ)}]{\gu{ગુણવું} (guṇvũ) — to multiply (numbers)}

\label{lesson:GU-C426-gunvu}



\begin{quote}
Before the new one: say the Gujarati for a crore, then the Gujarati for subtraction.
\end{quote}

\subsection*{You'll want to know: \gu{ગુણવું}}
\textbf{\gu{ગુણવું}} — \emph{guṇvũ} — \textquotedblleft{}to multiply (numbers)\textquotedblright{}.

\begin{quote}
\gu{મીરાએ} \gu{કાગળ} \gu{પર} \gu{બે} \gu{મોટી} \gu{રકમ} \gu{ગુણી}.
\end{quote}

\begin{quote}
\emph{Mīrāe kāgaḷ par be moṭī rakam guṇī.}
\end{quote}

\begin{quote}
Mira multiplied two large numbers on paper.
\end{quote}

Past, with \gu{તેણે} (he or she …): \gu{તેણે} \gu{રકમ} \gu{ગુણી}, \emph{teṇe rakam guṇī}.

\begin{cousinweb}
\textbf{\gu{ગુણવું}} (\emph{guṇvũ}) comes from Sanskrit \emph{guṇayati}, \textquotedblleft{}multiplies\textquotedblright{} (Monier-Williams), built on \emph{guṇa}, a strand or a fold: a number folded over on itself.
\end{cousinweb}

\begin{grammarlens}[title={What you've built}]
One more word for getting around.
\end{grammarlens}

\subsection*{Your turn}
\begin{itemize}
\raggedright
  \item \emph{Say it:} \emph{guṇvũ}
  \item \emph{Say it:} \emph{guṇvũ}, once more
  \item \emph{Say it:} \emph{Mīrāe kāgaḷ par be moṭī rakam guṇī.} — Mira multiplied two large numbers on paper.
  \item \emph{Recall:} say the Gujarati for a crore, then the Gujarati for subtraction, then say \emph{guṇvũ} again
\end{itemize}

\begin{tcolorbox}[breakable,colback=teal!4,colframe=teal!35!black,title={Before you move on}]
What does \gu{ગુણવું} mean? (\textquotedblleft{}To multiply (numbers)\textquotedblright{}.) Say it once more.
\end{tcolorbox}
\section[\texorpdfstring{\gu{વધારવું} (vadhārvũ)}{વધારવું (vadhārvũ)}]{\gu{વધારવું} (vadhārvũ) — to increase; to raise (a price)}

\label{lesson:GU-C426-vadharvu}



\begin{quote}
Before the new one: say the Gujarati for subtraction, then the Gujarati for to multiply.
\end{quote}

\subsection*{You'll want to know: \gu{વધારવું}}
\textbf{\gu{વધારવું}} — \emph{vadhārvũ} — \textquotedblleft{}to increase; to raise (a price)\textquotedblright{}.

\begin{quote}
\gu{દુકાનદારે} \gu{ભાવ} \gu{વધાર્યો}.
\end{quote}

\begin{quote}
\emph{dukāndāre bhāv vadhāryo.}
\end{quote}

\begin{quote}
The shopkeeper put the price up.
\end{quote}

Past, with \gu{તેણે} (he or she …): \gu{તેણે} \gu{ભાડું} \gu{વધાર્યું}, \emph{teṇe bhāḍũ vadhāryũ}.

\begin{cousinweb}
\textbf{\gu{વધારવું}} (\emph{vadhārvũ}) is built on \gu{વધારે}, more, a tadbhava of Sanskrit \emph{vṛddhatara}, \textquotedblleft{}larger\textquotedblright{} (CDIAL 12074): to make something more. \gu{વધવું}, to go up, is its partner.
\end{cousinweb}

\begin{grammarlens}[title={What you've built}]
One more word for getting around.
\end{grammarlens}

\subsection*{Your turn}
\begin{itemize}
\raggedright
  \item \emph{Say it:} \emph{vadhārvũ}
  \item \emph{Say it:} \emph{vadhārvũ}, once more
  \item \emph{Say it:} \emph{dukāndāre bhāv vadhāryo.} — The shopkeeper put the price up.
  \item \emph{Recall:} say the Gujarati for subtraction, then the Gujarati for to multiply, then say \emph{vadhārvũ} again
\end{itemize}

\begin{tcolorbox}[breakable,colback=teal!4,colframe=teal!35!black,title={Before you move on}]
What does \gu{વધારવું} mean? (\textquotedblleft{}To increase; to raise (a price)\textquotedblright{}.) Say it once more.
\end{tcolorbox}
